% ============================================================
%  Formalizing Quantitative Equational Reasoning
%  Program Verification Exam -- 07/09/2026
%
%  Compile with:  latexmk -pdf presentation.tex
% ============================================================
\documentclass[10pt,aspectratio=169]{beamer}

\usetheme[progressbar=frametitle,block=fill]{metropolis}

\usepackage[T1]{fontenc}
\usepackage[utf8]{inputenc}
\usepackage{amsmath,amssymb,amsfonts}
\usepackage{bm}
\usepackage{stmaryrd}          % \llbracket \rrbracket
\usepackage{mathtools}
\usepackage{mathpartir}        % inference rules
\usepackage{listings}
\usepackage{xcolor}
\usepackage[most]{tcolorbox}
\usepackage{tikz}
\usetikzlibrary{arrows.meta,positioning,decorations.pathmorphing}
\usepackage{booktabs}
\usepackage{nicematrix}           % row/column-labelled context matrices
\usepackage{appendixnumberbeamer}  % \inserttotalframenumber excludes the appendix

% ============================================================
%  Listings setup for Rocq source and for Rocq goal states
% ============================================================

\definecolor{rocqkw}{HTML}{5B4A7A}     % declaration heads only
\definecolor{rocqcon}{HTML}{B3541E}    % constructors / our names
\definecolor{rocqcom}{HTML}{6A737D}    % comments
\definecolor{goalbg}{HTML}{F4F6F8}
\definecolor{goalrule}{HTML}{C4CCD4}
\definecolor{srcbg}{HTML}{FBF9F4}

\lstdefinelanguage{Rocq}{
  % Only the declaration heads are marked; everything else stays plain.
  morekeywords=[1]{Record, Inductive, Definition, Fixpoint, Lemma, Theorem,
    Proposition, Corollary, Notation, Section, End, Variable, Hypothesis,
    Context, Proof, Qed, Defined},
  sensitive=true,
  morecomment=[n]{(*}{*)},
  morestring=[b]",
}

\lstdefinestyle{rocqsrc}{
  language=Rocq,
  basicstyle=\ttfamily\fontsize{6.6pt}{7.5pt}\selectfont,
  keywordstyle=[1]\color{rocqkw},
  commentstyle=\color{rocqcom}\itshape,
  showstringspaces=false,
  columns=fullflexible,
  keepspaces=true,
  breaklines=true,
  breakatwhitespace=false,
  literate={->}{{$\rightarrow$}}2 {=>}{{$\Rightarrow$}}2
           {<=}{{$\leq$}}2 {>=}{{$\geq$}}2 {|-}{{$\vdash$}}2
           {\\/}{{$\vee$}}2 {/\\}{{$\wedge$}}2,
  aboveskip=0.2em, belowskip=0.2em,
}

\lstdefinestyle{rocqgoal}{
  style=rocqsrc,
  basicstyle=\ttfamily\fontsize{6pt}{6.9pt}\selectfont,
}

\lstnewenvironment{rocq}[1][]{\lstset{style=rocqsrc,#1}}{}
\lstnewenvironment{rocqg}[1][]{\lstset{style=rocqgoal,#1}}{}

% A framed box for a proof-state dump
% A proof state, split so the conclusion is unmistakable:
%   upper part = hypotheses (context), lower part = what must be proved.
\definecolor{goalconcbg}{HTML}{E3EDF6}
\definecolor{goalconcfg}{HTML}{0C3149}

\lstnewenvironment{gctx}[1][]{\lstset{style=rocqgoal,#1}}{}
\lstnewenvironment{gconc}[1][]{%
  \lstset{style=rocqgoal,
    basicstyle=\ttfamily\fontsize{6.4pt}{7.3pt}\selectfont\color{goalconcfg},#1}}{}

\newtcolorbox{goalbox}[1][]{
  bicolor,
  colback=goalbg, colbacklower=goalconcbg,
  colframe=goalrule,
  boxrule=0.4pt, arc=1pt,
  left=2pt, right=2pt, top=1.5pt, bottom=1.5pt,
  middle=1.5pt,
  before skip=3pt, after skip=3pt,
  fonttitle=\scriptsize\bfseries,
  #1
}

\newtcblisting{srcbox}[1][]{
  listing only,
  listing options={style=rocqsrc},
  colback=srcbg, colframe=goalrule,
  boxrule=0.4pt, arc=1pt,
  left=2pt, right=2pt, top=1.5pt, bottom=1.5pt,
  before skip=3pt, after skip=3pt,
  fonttitle=\scriptsize\bfseries,
  #1
}

% ============================================================
%  Stepped proof scripts: a script pane whose "current" line is
%  highlighted, earlier lines normal, later lines dimmed.
%
%    \pstep{3}                 -- declare which step this slide shows
%    \pline{0}{...}            -- context line, never dimmed
%    \pline{n}{...}            -- line belonging to step n
% ============================================================

\definecolor{stepbg}{HTML}{FBEDC6}     % the line you are on
\definecolor{stepdim}{HTML}{A9B1B9}    % not executed yet

\newcommand{\pcur}{0}
\newcommand{\pstep}[1]{\renewcommand{\pcur}{#1}}
% a strut exactly one \baselineskip tall (so it tracks the pane's font
% size), making consecutive highlighted lines merge into one block
\newcommand{\pstrut}{\vrule height 0.76\baselineskip depth 0.24\baselineskip width 0pt}
\newcommand{\pline}[2]{%
  \ifnum#1=0\relax
    \makebox[\linewidth][l]{\pstrut#2}%
  \else\ifnum#1<\pcur\relax
    \makebox[\linewidth][l]{\pstrut#2}%
  \else\ifnum#1=\pcur\relax
    {\setlength{\fboxsep}{0pt}%
     \colorbox{stepbg}{\makebox[\linewidth][l]{\pstrut#2}}}%
  \else
    \makebox[\linewidth][l]{\pstrut\color{stepdim}#2}%
  \fi\fi\fi\par}

\newtcolorbox{scriptbox}[1][]{
  colback=srcbg, colframe=goalrule,
  boxrule=0.4pt, arc=1pt,
  left=2pt, right=2pt, top=1.5pt, bottom=1.5pt,
  before skip=3pt, after skip=3pt,
  fontupper=\ttfamily\fontsize{6pt}{7.6pt}\selectfont\frenchspacing,
  #1
}


% The proof script of [frel_derives_weakening], shared by its walkthrough
% slides; \pstep{n} on each slide selects the line that is current.
% The proof script of [derives_fin_sound]; the twelve induction cases are
% shown one line each, tagged with the rule they come from.
\newcommand{\ptag}[1]{\hfill\textit{#1}}
% The proof script of [kantorovich_models_ica] (Proposition 4.6).
\newcommand{\modelscript}{%
\pline{0}{Theorem kantorovich\_models\_ica :}
\pline{0}{\ \ fuzzy\_models kantorovich\_algebra (@ica\_theory R).}
\pline{1}{Proof.}
\pline{1}{\ \ rewrite /fuzzy\_models.}
\pline{1}{\ \ move => Y phi Hphi.}
\pline{1}{\ \ case: Hphi => [p | p | p q | p eps delta ...] rho.}
\pline{2}{\ \ - by rewrite kantorovich\_eval\_op}
\pline{2}{\ \ \ \ \ \ convex\_mixture\_idem.\ptag{Idem}}
\pline{2}{\ \ - by rewrite !kantorovich\_eval\_op}
\pline{2}{\ \ \ \ \ \ convex\_mixture\_skew\_comm.\ptag{SkewComm}}
\pline{2}{\ \ - by rewrite !kantorovich\_eval\_op}
\pline{2}{\ \ \ \ \ \ convex\_mixture\_skew\_assoc.\ptag{SkewAssoc}}
\pline{3}{\ \ - rewrite !kantorovich\_eval\_op.\ptag{Interp}}
\pline{3}{\ \ \ \ apply: (le\_trans}
\pline{3}{\ \ \ \ \ \ (kantorovich\_convex\_mixture \_ \_ \_ \_ \_ Hd)).}
\pline{3}{\ \ \ \ apply: weighted\_sum\_le.}
\pline{4}{\ \ \ \ - apply: (le\_trans (interpretation\_nexp}
\pline{4}{\ \ \ \ \ \ \ \ rho (inord 0) (inord 2))).}
\pline{4}{\ \ \ \ \ \ by rewrite /= /ica\_interp\_rel !inordK.}
\pline{4}{\ \ \ \ - apply: (le\_trans (interpretation\_nexp}
\pline{4}{\ \ \ \ \ \ \ \ rho (inord 1) (inord 3))).}
\pline{4}{\ \ \ \ \ \ by rewrite /= /ica\_interp\_rel !inordK.}
\pline{4}{Qed.}
}

\newcommand{\soundnessscript}{%
\pline{0}{Theorem derives\_fin\_sound (A : FuzzyAlgebra R sig)}
\pline{0}{\ \ \ \ (E : fuzzy\_theory R sig) (X : fuzzy\_space R)}
\pline{0}{\ \ \ \ (phi : @judgement R sig (fcarrier X)) :}
\pline{0}{\ \ fuzzy\_models A E -> derives\_fin E phi ->}
\pline{0}{\ \ satisfies A X phi.}
\pline{1}{Proof.}
\pline{2}{\ \ rewrite /fuzzy\_models /satisfies /derives\_fin.}
\pline{2}{\ \ move => Hmodel Hderive rho.}
\pline{3}{\ \ remember FRelFinite as mode eqn:Hmode in Hderive.}
\pline{3}{\ \ revert Hmode.}
\pline{4}{\ \ induction Hderive; intros Hmode.}
\pline{5}{\ \ - apply: Hmodel X phi H rho.\ptag{Init}}
\pline{5}{\ \ - by [].\ptag{EqRefl}}
\pline{5}{\ \ - symmetry. exact: IHHderive ...\ptag{EqSym}}
\pline{5}{\ \ - by rewrite (IHHderive1 ...) ...\ptag{EqTrans}}
\pline{5}{\ \ - congr fa\_ops; apply ...\ptag{EqCong}}
\pline{6}{\ \ - apply: (interpretation\_nexp rho).\ptag{UseVars}}
\pline{7}{\ \ - apply: satisfies\_subst. ...\ptag{Subst}}
\pline{8}{\ \ - rewrite (IHHderive1 ...). ...\ptag{QEqRepL}}
\pline{8}{\ \ - rewrite -(IHHderive2 ...). ...\ptag{QEqRepR}}
\pline{8}{\ \ - apply: le\_trans ... H.\ptag{Up}}
\pline{8}{\ \ - apply: (andP (frel\_range ...)).2\ptag{Max}}
\pline{8}{\ \ - discriminate Hmode.\ptag{OrderCompl}}
\pline{8}{Qed.}
}

\newcommand{\weakeningscript}{%
\pline{0}{Lemma frel\_derives\_weakening}
\pline{0}{\ \ \ \ \{R : realType\} \{sig : signature\}}
\pline{0}{\ \ \ \ (E F : fuzzy\_theory R sig)}
\pline{0}{\ \ \ \ (mode : frel\_derivation\_mode)}
\pline{0}{\ \ \ \ (X : fuzzy\_space R)}
\pline{0}{\ \ \ \ (phi : @judgement R sig (fcarrier X)) :}
\pline{0}{\ \ fuzzy\_theory\_incl E F ->}
\pline{0}{\ \ frel\_derives E mode phi ->}
\pline{0}{\ \ frel\_derives F mode phi.}
\pline{1}{Proof.}
\pline{2}{\ \ rewrite /fuzzy\_theory\_incl.}
\pline{2}{\ \ move => Hincl H.}
\pline{3}{\ \ induction H.}
\pline{4}{\ \ - apply: FD\_Init.}
\pline{5}{\ \ \ \ apply: Hincl.}
\pline{6}{\ \ \ \ exact: H.}
\pline{7}{\ \ all: eauto using frel\_derives.}
\pline{7}{Qed.}
}

% ------------------------------------------------------------
%  Notation
% ------------------------------------------------------------
\NiceMatrixOptions{code-for-first-row=\scriptstyle,code-for-first-col=\scriptstyle}
% \ctx{<col spec>}{<body: first row = column labels, first col = row labels>}
\newcommand{\ctx}[2]{%
  \left(\;\begin{bNiceArray}{#1}[first-row,first-col,small]#2\end{bNiceArray}\right)}
\newcommand{\qeq}[1]{=_{#1}}
\newcommand{\den}[1]{\llbracket #1 \rrbracket}
\newcommand{\ICA}{\mathbb{I}\mathbb{C}\mathbb{A}}
\newcommand{\Dist}{\mathcal{D}}
\newcommand{\Terms}{\mathbb{T}}
\newcommand{\eps}{\varepsilon}
\newcommand{\Kan}{\mathcal{K}}
\newcommand{\Cpl}{\Gamma}
\newcommand{\bpl}[1]{+_{#1}}
\newcommand{\file}[1]{{\color{rocqcon}\texttt{#1}}}
\newcommand{\rid}[1]{{\ttfamily\small #1}}

% metropolis: make frametitles and the title page a bit tighter
\metroset{sectionpage=progressbar,numbering=fraction}
\makeatletter
\setlength{\metropolis@frametitle@padding}{1.1ex}
\makeatother
\setbeamerfont{title}{size=\Large,series=\bfseries}
\setbeamerfont{author}{size=\normalsize}
\setbeamerfont{institute}{size=\small}
\setbeamerfont{date}{size=\small}
\setbeamerfont{frametitle}{size=\normalsize,series=\bfseries}
\setbeamersize{text margin left=5mm,text margin right=5mm}

\title{Formalizing Quantitative Equational Reasoning}
\author{Benjamin Bennetzen}
\institute{Aalborg University \quad\textbullet\quad \texttt{bcb@aau.dk}\\[0.3em]
  {\small Supervised by Giorgio Bacci and Robin Kaarsgaard}}
\date{Program Verification Exam \quad\textbullet\quad 07/09/2026}

\begin{document}

\begin{frame}[plain]
  \titlepage
\end{frame}

% ============================================================
\section{Setting the scene}
% ============================================================

\begin{frame}{Motivation}
  \begin{itemize}
    \item[--] Equations are fundamental to mathematical reasoning.
      \medskip
    \item[--] Reasoning about programs is also based on program
      equivalences.
      \medskip
    \item[--] We want quantitative analogues of equational
      reasoning.
  \end{itemize}
\end{frame}

\begin{frame}{Replace standard equations with quantitative ones}
  \begin{center}
    \begin{tcolorbox}[colback=srcbg,colframe=goalrule,width=0.55\textwidth,
                      halign=center,boxrule=0.6pt,arc=8pt]
      \Huge $s \qeq{\eps} t$
    \end{tcolorbox}
  \end{center}
  \pause
  \vfill
  \begin{center}
    \large ``$s$ is within $\eps$ distance of $t$''
  \end{center}
\end{frame}

% ============================================================
\section{The classical picture: barycentric algebras}
% ============================================================

\begin{frame}{An equational theory of barycentric algebras}
  \begin{columns}[T]
    \begin{column}{0.42\textwidth}
      \textbf{Variables}\\[0.3em]
      $s,t,u,\dots$
    \end{column}
    \begin{column}{0.55\textwidth}
      \textbf{Operations}\\[0.3em]
      $\Sigma_{\mathrm{CA}} = \{\,\bpl{p} \mid 0 < p < 1\,\}$,\\
      each $\bpl{p}$ binary.
    \end{column}
  \end{columns}
  \vfill
  \textbf{Equations}
  \begin{align*}
    \text{(Idempotency)}\qquad & s \bpl{p} s = s\\[0.3em]
    \text{(Skew commutativity)}\qquad & s \bpl{p} t = t \bpl{1-p} s\\[0.3em]
    \text{(Skew associativity)}\qquad & (s \bpl{p} t) \bpl{q} u
      = s \bpl{pq} \bigl(t \bpl{\frac{(1-p)q}{1-pq}} u\bigr)
  \end{align*}
  \vfill
  {\footnotesize Also known as \emph{convex algebras} or \emph{convex spaces}.}
\end{frame}

\begin{frame}{A model of the theory}
  A model interprets each $\bpl{p}$ while respecting the axioms.
  \vfill
  One model: the finitely supported probability distributions
  $\Dist(X)$ with convex combinations
  \[
    (\mu \bpl{p} \nu)(x) \;=\; p\,\mu(x) + (1-p)\,\nu(x).
  \]
  \vfill
  \pause
  \textbf{Small example.} $X = \{a,b,c\}$, written as vectors
  $(\mu(a),\mu(b),\mu(c))$:
  \[
    \delta_a \bpl{\frac12} \delta_b = (\tfrac12,\tfrac12,0),
    \qquad
    \delta_a \bpl{p} \delta_a = \delta_a .
  \]
  \pause
  Skew associativity with $p = q = \tfrac12$
  (so $pq = \tfrac14$ and $\tfrac{(1-p)q}{1-pq} = \tfrac13$):
  \[
    \underbrace{(\delta_a \bpl{\frac12} \delta_b) \bpl{\frac12} \delta_c}
      _{(\frac14,\frac14,\frac12)}
    \;=\;
    \underbrace{\delta_a \bpl{\frac14}
      (\delta_b \bpl{\frac13} \delta_c)}_{(\frac14,\frac14,\frac12)}
  \]
  \vfill
  {\footnotesize In fact $\Dist(X)$ is the \emph{free} convex algebra on $X$.}
\end{frame}

% ============================================================
\section{Formalizing the syntax}
% ============================================================

\begin{frame}[fragile]{Formalization: signatures and terms}
  \file{src/Signature.v}
  \begin{srcbox}
Record signature := {
  sym : Type;
  arity : sym -> nat;
}.

Inductive term (sig : signature) (X : Type) : Type :=
  | Var : X -> term sig X
  | App : forall f : sym sig, ('I_(arity f) -> term sig X) -> term sig X.

Definition algebra_ops (sig : signature) (A : Type) : Type :=
  forall f : sym sig, (ordinal (arity f) -> A) -> A.
  \end{srcbox}
  \begin{itemize}
    \item \rid{sym} is an arbitrary \rid{Type}: the signature may be
      uncountable, as $\Sigma_{\mathrm{CA}}$ is.
    \item Arguments are given as functions out of the ordinal type
      \rid{'I\_n}, so \rid{term} stays strictly positive and MathComp's
      \rid{fintype} machinery applies.
  \end{itemize}
\end{frame}

\begin{frame}[fragile]{Formalization: substitution}
  \begin{srcbox}
Fixpoint subst_term {sig X Y} (sigma : X -> term sig Y)
    (t : term sig X) : term sig Y :=
  match t with
  | Var x => sigma x
  | App f args => App f (subst_term sigma \o args)
  end.

Lemma subst_term_comp {sig X Y Z}
    (sigma : Y -> term sig Z) (tau : X -> term sig Y) (t : term sig X) :
  subst_term sigma (subst_term tau t) = subst_term (subst_comp sigma tau) t.

Lemma subst_term_var {sig X} (t : term sig X) :
  subst_term Var t = t.
  \end{srcbox}
  \vfill
  Terms form a monad on \rid{Type}; these two lemmas are its laws.
  Both go through by induction plus
  \rid{functional\_extensionality} on the argument vectors.
\end{frame}

% ============================================================
\section{Going quantitative}
% ============================================================

\begin{frame}{How to extend to the quantitative setting}
  \begin{description}
    \item[\underline{Quantitative} equational theories]\hfill\\
      operations $+$ \underline{quantitative} equations.
      \medskip
    \item[\underline{EMet}-Lawvere theories]\hfill\\
      Lawvere categories \underline{enriched in the category of extended
      metric spaces}.
      \medskip
    \item[Monads on $\mathbf{Set}$ (\underline{EMet})]\hfill\\
      \underline{quantitative} effects.
  \end{description}
  \vfill
  \hfill{\small Mardare, Panangaden and Plotkin (LICS'16, LICS'17);}\\
  \hfill{\small Bacci, Mardare, Panangaden and Plotkin.}
\end{frame}

\begin{frame}{What is a model in the quantitative setting?}
  If
  \[ s \qeq{\eps} t \]
  then for an interpretation $\den{\cdot} : \Terms_\Sigma(\mathrm{Var}) \to A$
  \[ d\bigl(\den{s},\den{t}\bigr) \;\leq\; \eps \]
  where $d$ is a metric.
  \vfill
  \pause
  \begin{alertblock}{Note}
    $s \qeq{0} t$ means strict equality, because $d$ is a metric:
    $d(a,b) = 0 \Rightarrow a = b$.
  \end{alertblock}
\end{frame}

\begin{frame}{Relaxing to fuzzy relations}
  A generalization relaxes the metric requirement to simply a
  fuzzy relation
  \[ d : A \times A \longrightarrow [0,1] . \]
  \vfill
  \textbf{No additional assumptions:} not symmetric, no triangle
  inequality, and $d(a,a)$ need not be $0$.
  \vfill
  \pause
  \begin{alertblock}{Consequence}
    We still need $=$ as a separate judgement form denoting exact
    equality, since $s \qeq{0} t$ may \emph{not} imply
    $\den{s} = \den{t}$.
  \end{alertblock}
  \vfill
  \hfill{\small Mio, Sarkis and Vignudelli, LMCS 20 (2024).}
\end{frame}

\begin{frame}{Satisfiability relation}
  Let $\mathbb{A} = (A, \{op^{\mathbb{A}}\}, d_{\mathbb{A}})$ be a
  quantitative algebra.
  \vfill
  \[
    \mathbb{A} \models \forall (B,d_B).\, s = t
    \quad\text{iff}\quad
    \den{s} = \den{t}
  \]
  \vspace{-1em}
  \[
    \mathbb{A} \models \forall (B,d_B).\, s \qeq{\eps} t
    \quad\text{iff}\quad
    d_{\mathbb{A}}\bigl(\den{s},\den{t}\bigr) \leq \eps
  \]
  \vfill
  for all $1$-Lipschitz $\den{\cdot} : (B,d_B) \to
  (A,d_\mathbb{A})$.
  \vfill
  \pause
  {\small The context $(B,d_B)$ is itself a fuzzy relation on the
  variables. Taking $d_B \equiv 1$ (the \emph{discrete} fuzzy relation)
  makes every interpretation $1$-Lipschitz, recovering ordinary
  universal quantification over variables.}
\end{frame}

% ============================================================
\section{Formalizing quantitative algebra}
% ============================================================

\begin{frame}[fragile]{Formalization: fuzzy spaces and algebras}
  \file{src/FuzzyRelation.v}
  \begin{srcbox}
Record fuzzy_space (R : realType) := {
    fcarrier : Type;
    frel : fcarrier -> fcarrier -> R;
    frel_range : forall x y, (0 <= frel x y <= 1)%R;
}.

Definition frel_nonexpansive {R : realType} (A B : fuzzy_space R)
    (f : fcarrier A -> fcarrier B) : Prop :=
  forall x y, (frel B (f x) (f y) <= frel A x y)%R.

Record FuzzyAlgebra (R : realType) (sig : signature) := {
  fa_space : fuzzy_space R;
  fa_ops : algebra_ops sig (fcarrier fa_space)
}.
  \end{srcbox}
  \begin{itemize}
    \item \rid{R : realType} is MathComp-Analysis' abstract real type.
    \item A \rid{FuzzyAlgebra} carries no non-expansiveness
      obligation on \rid{fa\_ops} --- that is expressed by
      quantitative equations instead.
  \end{itemize}
\end{frame}

\begin{frame}[fragile]{Formalization: judgements and interpretations}
  \begin{columns}[T]
  \begin{column}{0.5\textwidth}
  \begin{srcbox}
Inductive judgement {R : realType}
    (sig : signature) (X : Type) :=
  | EqJ  : term sig X -> term sig X ->
           judgement sig X
  | QEqJ : R -> term sig X -> term sig X ->
           judgement sig X.
  \end{srcbox}
  \end{column}
  \begin{column}{0.5\textwidth}
  \begin{srcbox}
Record interpretation
    {R : realType} {sig}
    (X : fuzzy_space R)
    (A : FuzzyAlgebra R sig) := {
  interpretation_fun :>
    fcarrier X -> fcarrier (fa_space A);
  interpretation_nexp :
    frel_nonexpansive X (fa_space A)
      interpretation_fun;
}.
  \end{srcbox}
  \end{column}
  \end{columns}
  \vfill
  \begin{srcbox}
Fixpoint fuzzy_eval {R : realType} {sig X} (A : FuzzyAlgebra R sig)
    (rho : X -> fcarrier (fa_space A)) (t : term sig X) : fcarrier (fa_space A) :=
  match t with
  | Var x => rho x
  | App f args => fa_ops A f (fun i => fuzzy_eval A rho (args i))
  end.
  \end{srcbox}
  The $1$-Lipschitz side condition is packed into the record: a
  value of type \rid{interpretation X A} \emph{is} an admissible
  $\den{\cdot}$.
\end{frame}

\begin{frame}[fragile]{Formalization: satisfaction, theories, models}
  \begin{srcbox}
Definition satisfies {R : realType} {sig} (A : FuzzyAlgebra R sig)
    (X : fuzzy_space R) (phi : judgement sig (fcarrier X)) : Prop :=
  forall rho : interpretation X A,
    match phi with
    | EqJ s t => fuzzy_eval A rho s = fuzzy_eval A rho t
    | QEqJ eps s t =>
      (frel (fa_space A) (fuzzy_eval A rho s) (fuzzy_eval A rho t) <= eps)%R
    end.

Definition fuzzy_theory (R : realType) (sig : signature) :=
  forall X : fuzzy_space R, @judgement R sig (fcarrier X) -> Prop.

Definition fuzzy_models {R : realType} {sig} (A : FuzzyAlgebra R sig)
    (E : fuzzy_theory R sig) : Prop :=
  forall X phi, E X phi -> satisfies A X phi.
  \end{srcbox}
  \vfill
  This is exactly the satisfiability relation from the previous slide:
  \rid{satisfies} is $\models$ for a single judgement, and
  \rid{fuzzy\_models} is $\mathbb{A} \models E$.
  \vfill
  A \rid{fuzzy\_theory} is a predicate over all contexts at
  once --- necessary, since axiom schemes such as
  ``interpolative'' range over every fuzzy relation.
\end{frame}

% ============================================================
\section{The quantitative theory $\ICA$}
% ============================================================

\begin{frame}{A quantitative theory of barycentric algebras}
  \begin{align*}
    \text{(Idempotency)}\quad &
      \forall\ctx{c}{ & x \\ x & 1 }.\;
      x \bpl{p} x = x
    \\[0.4em]
    \text{(Skew comm.)}\quad &
      \forall\ctx{cc}{ & x & y \\ x & 1 & 1 \\ y & 1 & 1 }.\;
      x \bpl{p} y = y \bpl{1-p} x
    \\[0.4em]
    \text{(Skew assoc.)}\quad &
      \forall\ctx{ccc}{ & x & y & z \\
        x & 1 & 1 & 1 \\ y & 1 & 1 & 1 \\ z & 1 & 1 & 1 }.\;
      (x \bpl{p} y) \bpl{q} z = x \bpl{pq}
        \bigl(y \bpl{\frac{(1-p)q}{1-pq}} z\bigr)
    \\[0.6em]
    \text{(Interpolative)}\quad &
      \forall\ctx{cccc}{ & x & y & w & z \\
        x & 1 & 1 & \eps & 1 \\
        y & 1 & 1 & 1 & \delta \\
        w & 1 & 1 & 1 & 1 \\
        z & 1 & 1 & 1 & 1 }.\;
      x \bpl{p} y \;\qeq{p\eps + (1-p)\delta}\; w \bpl{p} z
  \end{align*}
  \vfill
  The matrix is the fuzzy relation $d_B$ on the variables: it
  constrains the admissible interpretations to those with
  $d(\den{x},\den{w}) \leq \eps$ and $d(\den{y},\den{z}) \leq \delta$.
  \hfill{\footnotesize (the theory $\ICA$)}
\end{frame}

\begin{frame}{An aside about moving piles of sand}
  \begin{center}
    \begin{tikzpicture}[scale=0.9]
      % left pile: two bumps
      \draw[thick,rounded corners=1pt]
        (0,0) .. controls (0.4,0.1) and (0.5,1.4) .. (0.9,1.4)
              .. controls (1.2,1.4) and (1.2,0.5) .. (1.45,0.5)
              .. controls (1.7,0.5) and (1.75,1.5) .. (2.1,1.5)
              .. controls (2.5,1.5) and (2.6,0.1) .. (3.1,0);
      % right pile: one bump
      \draw[thick,rounded corners=1pt]
        (6.2,0) .. controls (6.9,0.1) and (7.1,1.6) .. (7.8,1.6)
              .. controls (8.5,1.6) and (8.7,0.1) .. (9.4,0);
      % figure with shovel
      \begin{scope}[shift={(3.9,0.0)},scale=0.9]
        \draw[thick] (0.5,1.9) circle (0.28);
        \fill (0.42,1.95) circle (0.03); \fill (0.58,1.95) circle (0.03);
        \draw[thick] (0.5,1.62) -- (0.5,0.85);
        \draw[thick] (0.5,1.45) -- (0.05,1.05);
        \draw[thick] (0.5,1.45) -- (0.95,1.15);
        \draw[thick] (0.5,0.85) -- (0.2,0.05);
        \draw[thick] (0.5,0.85) -- (0.8,0.05);
        \draw[thick] (0.05,1.05) -- (-0.25,0.35);
        \draw[thick,fill=goalbg] (-0.42,0.35) -- (-0.08,0.35)
          -- (-0.13,0.05) -- (-0.37,0.05) -- cycle;
      \end{scope}
      \draw[-{Stealth[length=3mm]},thick] (0.3,-0.45) -- (9.1,-0.45);
    \end{tikzpicture}
  \end{center}
  \vfill
  Out of all the possible ways to move the sand, which is the least
  amount of work?
  \[
    \Kan(d)(\mu,\nu) \;=\; \inf_{\gamma \in \Cpl(\mu,\nu)}
      \sum_{a,b \in A} \gamma(a,b) \cdot d(a,b)
    \qquad
    \substack{\text{\emph{optimal}}\\ \text{\emph{transport}}}
  \]
  \vfill
  {\footnotesize $\Cpl(\mu,\nu)$: the joint distributions with marginals
  $\mu$ and $\nu$. Always non-empty --- take the independent product.}
\end{frame}

\begin{frame}{A model in the quantitative setting}
  The finitely supported probability distributions, their convex sums,
  and the Kantorovich ``lifting'' form a model.
  \vfill
  \pause
  Actually, it is axiomatised by $\ICA$:
  \vfill
  \begin{block}{Soundness}
    If $s \qeq{\eps} t$ is derivable, then
    $\Kan(d)(\den{s},\den{t}) \leq \eps$.
  \end{block}
  \vfill
  \begin{block}{Completeness}
    If $\Kan(d)(\den{s},\den{t}) \leq \eps$, then $s \qeq{\eps} t$ is
    derivable.
  \end{block}
\end{frame}

% ============================================================
\section{The deduction system}
% ============================================================

\begin{frame}{General axioms}
  Some things should always hold:
  \[
    \inferrule{a = b \\ b = c}{a = c}
  \]
  \vfill
  \pause
  Does not quite hold for $\qeq{\eps}$:
  \begin{columns}[c]
    \begin{column}{0.45\textwidth}
      \[
        \inferrule{a \qeq{\eps} b \\ b \qeq{\delta} c}
                  {a \qeq{\eps + \delta} c}
      \]
    \end{column}
    \begin{column}{0.5\textwidth}
      \centering
      \begin{tikzpicture}[scale=1.1]
        \node (a) at (0,0) {$a$};
        \node (c) at (2.2,0) {$c$};
        \node (b) at (1.1,1.1) {$b$};
        \draw (a) -- node[above left,font=\small] {$\eps$} (b);
        \draw (b) -- node[above right,font=\small] {$\delta$} (c);
        \draw (a) -- node[below,font=\small] {$\leq \eps + \delta$} (c);
      \end{tikzpicture}
    \end{column}
  \end{columns}
  \vfill
  \pause
  \begin{alertblock}{}
    \dots this is only general in the metric view. For an
    arbitrary fuzzy relation there is no triangle inequality, so the
    rule is unsound.
  \end{alertblock}
\end{frame}

\begin{frame}{Other general axioms}
  \begin{columns}[c]
    \begin{column}{0.5\textwidth}
      \[
        \inferrule{\eps \leq \delta \\ a \qeq{\eps} b}{a \qeq{\delta} b}
      \]
    \end{column}
    \begin{column}{0.4\textwidth}
      \[
        \inferrule{ }{a \qeq{1} b}
      \]
    \end{column}
  \end{columns}
  \vfill\vfill
  \[
    \inferrule{\{\, s \qeq{\delta} t \mid \eps < \delta \,\}}
              {s \qeq{\eps} t}
  \]
  \vfill
  \begin{itemize}
    \item[$\hookrightarrow$] This rule has infinite width: proofs
      are no longer finite objects.
    \item[$\hookrightarrow$] Known as the \emph{order-completeness},
      \emph{Archimedean}, or \emph{continuity} rule.
    \item[$\hookrightarrow$] \alert{Proven to be not needed} for
      barycentric algebras: $\ICA$ is a \emph{compact} theory.
  \end{itemize}
  \vfill
  \hfill{\small Mio, \emph{Compact Quantitative Theories of Convex
  Algebras}, MFPS 2025.}
\end{frame}

\begin{frame}[fragile]{Formalization: the metric-view system}
  \file{src/QET.v} --- the Mardare--Panangaden--Plotkin system, over
  extended metrics, with conditional inferences $\Gamma \vdash \phi$:
  \begin{srcbox}
Inductive derives {R : realType} {sig} (U : axiom_scheme R sig)
    : forall X : Type, ctx R sig X -> qeq R sig X -> Prop :=
  | D_Refl   : forall X Gamma t, derives U Gamma (t ~[0] t)
  | D_Symm   : forall X Gamma t s eps, Rnn eps ->
      derives U Gamma (t ~[eps] s) -> derives U Gamma (s ~[eps] t)
  | D_Triang : forall X Gamma t s u eps eps', Rnn eps -> Rnn eps' ->
      derives U Gamma (t ~[eps] s) -> derives U Gamma (s ~[eps'] u) ->
      derives U Gamma (t ~[eps + eps'] u)
  | D_Arch   : forall X Gamma t s eps, Rnn eps ->
      (forall eps', (eps < eps')%R -> derives U Gamma (t ~[eps'] s)) ->
      derives U Gamma (t ~[eps] s)
  | D_NExp | D_Subst | D_Cut | D_Assumpt | D_Axiom (* ... *).
  \end{srcbox}
  \vfill
  \rid{D\_Arch} is the infinitary rule: its premise is a
  function over all $\eps' > \eps$, so a derivation tree is
  well-founded but of infinite width.
\end{frame}

\begin{frame}[fragile]{Formalization: the fuzzy-relation system}
  \file{src/FRelDeduction.v} --- the strict-equality fragment is
  ordinary Birkhoff:
  \begin{srcbox}
Inductive frel_derives {R : realType} {sig : signature}
    (E : fuzzy_theory R sig) : frel_derivation_mode ->
    forall X : fuzzy_space R, @judgement R sig (fcarrier X) -> Prop :=
  | FD_Init : forall mode X phi,
      E X phi -> frel_derives E mode phi
  | FD_EqRefl : forall mode X (s : term sig (fcarrier X)),
      frel_derives E mode (EqJ s s)
  | FD_EqSym : forall mode X (s t : term sig (fcarrier X)),
      frel_derives E mode (EqJ s t) -> frel_derives E mode (EqJ t s)
  | FD_EqTrans : forall mode X (s t u : term sig (fcarrier X)),
      frel_derives E mode (EqJ s t) -> frel_derives E mode (EqJ t u) ->
      frel_derives E mode (EqJ s u)
  | FD_EqCong : forall mode X (f : sym sig)
      (s t : 'I_(arity f) -> term sig (fcarrier X)),
      (forall i, frel_derives E mode (EqJ (s i) (t i))) ->
      frel_derives E mode (EqJ (App f s) (App f t))
  (* ... *)
  \end{srcbox}
  Quantitative reflexivity, symmetry and transitivity are
  \alert{deliberately absent}: an arbitrary fuzzy relation need not
  satisfy them.
\end{frame}

\begin{frame}[fragile]{Formalization: the quantitative rules}
  \begin{srcbox}
  | FD_UseVariables : forall mode X (x y : fcarrier X),
      frel_derives E mode (QEqJ (frel X x y) (Var x) (Var y))
  | FD_Subst : forall mode X Y phi
      (sigma : fcarrier X -> term sig (fcarrier Y)),
      frel_derives E mode phi ->
      (forall x y, frel_derives E mode
          (QEqJ (frel X x y) (sigma x) (sigma y))) ->
      frel_derives E mode (subst_judgement sigma phi)
  | FD_QEqReplaceL | FD_QEqReplaceR (* rewriting with EqJ *)
  | FD_Up : forall mode X eps delta (s t : term sig (fcarrier X)),
      (eps <= delta)%R ->
      frel_derives E mode (QEqJ eps s t) -> frel_derives E mode (QEqJ delta s t)
  | FD_Max : forall mode X (s t : term sig (fcarrier X)),
      frel_derives E mode (QEqJ 1 s t)
  | FD_OrderComplete : forall X eps (s t : term sig (fcarrier X)),
      (forall delta, (eps < delta)%R ->
        frel_derives E FRelFull (QEqJ delta s t)) ->
      frel_derives E FRelFull (QEqJ eps s t).
  \end{srcbox}
  \begin{itemize}\small
    \item \rid{FD\_Subst}'s second premise is exactly ``$\sigma$ is
      \emph{provably} $1$-Lipschitz''.
    \item \rid{FD\_OrderComplete} only fires at \rid{FRelFull}, so
      \rid{derives\_fin} is the finitary sub-system, inside the
      same inductive type.
  \end{itemize}
\end{frame}

% ============================================================
\section{A lemma, in full}
% ============================================================

\begin{frame}[fragile]{Lemma \rid{frel\_derives\_weakening}}
  \begin{columns}[T]
  \begin{column}{0.42\textwidth}
    \begin{scriptbox}\pstep{1}\weakeningscript\end{scriptbox}
  \end{column}
  \begin{column}{0.56\textwidth}
    \begin{goalbox}[title={Proof state}]
\begin{gctx}
R : realType
sig : signature
E, F : fuzzy_theory R sig
mode : frel_derivation_mode
X : fuzzy_space R
phi : judgement sig (fcarrier X)
\end{gctx}
\tcblower
\begin{gconc}
|- fuzzy_theory_incl E F ->
   frel_derives E mode phi ->
   frel_derives F mode phi
\end{gconc}
    \end{goalbox}
    \smallskip
    {\footnotesize \emph{Enlarging the theory never loses a
    derivation.}\\[0.4em]
    In every proof state: the grey half is the \textbf{context}, the
    blue half is \textbf{what must be proved}.}
  \end{column}
  \end{columns}
\end{frame}

\begin{frame}[fragile]{Lemma \rid{frel\_derives\_weakening}}
  \begin{columns}[T]
  \begin{column}{0.42\textwidth}
    \begin{scriptbox}\pstep{2}\weakeningscript\end{scriptbox}
  \end{column}
  \begin{column}{0.56\textwidth}
    \begin{goalbox}[title={Proof state}]
\begin{gctx}
R : realType
sig : signature
E, F : fuzzy_theory R sig
mode : frel_derivation_mode
X : fuzzy_space R
phi : judgement sig (fcarrier X)
Hincl : forall (X0 : fuzzy_space R)
          (phi0 : judgement sig (fcarrier X0)),
        E X0 phi0 -> F X0 phi0
H : frel_derives E mode phi
\end{gctx}
\tcblower
\begin{gconc}
|- frel_derives F mode phi
\end{gconc}
    \end{goalbox}
    \smallskip
    {\footnotesize The hypothesis \rid{H} is the derivation we induct
    on. \rid{X} and \rid{phi} are indices of \rid{frel\_derives}, so
    \rid{induction} generalizes them automatically.}
  \end{column}
  \end{columns}
\end{frame}

\begin{frame}[fragile]{Lemma \rid{frel\_derives\_weakening}}
  \begin{columns}[T]
  \begin{column}{0.42\textwidth}
    \begin{scriptbox}\pstep{3}\weakeningscript\end{scriptbox}
    \smallskip
    {\footnotesize \textbf{12 goals} --- one per constructor of
    \rid{frel\_derives}. Only the first is interesting.}
  \end{column}
  \begin{column}{0.56\textwidth}
    \begin{goalbox}[title={Goal 1 / 12 \quad (FD\_Init)}]
\begin{gctx}
Hincl : forall (X0 : fuzzy_space R)
          (phi0 : judgement sig (fcarrier X0)),
        E X0 phi0 -> F X0 phi0
H : E X phi
\end{gctx}
\tcblower
\begin{gconc}
|- frel_derives F mode phi
\end{gconc}
    \end{goalbox}
    \begin{goalbox}[title={Goal 4 / 12 \quad (FD\_EqTrans), typical}]
\begin{gctx}
s, t, u : term sig (fcarrier X)
H  : frel_derives E mode (EqJ s t)
H0 : frel_derives E mode (EqJ t u)
IHfrel_derives1 : frel_derives F mode (EqJ s t)
IHfrel_derives2 : frel_derives F mode (EqJ t u)
\end{gctx}
\tcblower
\begin{gconc}
|- frel_derives F mode (EqJ s u)
\end{gconc}
    \end{goalbox}
  \end{column}
  \end{columns}
\end{frame}

\begin{frame}[fragile]{Lemma \rid{frel\_derives\_weakening}}
  \begin{columns}[T]
  \begin{column}{0.42\textwidth}
    \begin{scriptbox}\pstep{4}\weakeningscript\end{scriptbox}
  \end{column}
  \begin{column}{0.56\textwidth}
    \begin{goalbox}[title={Goal 1 / 12, after \texttt{apply: FD\_Init}}]
\begin{gctx}
Hincl : forall (X0 : fuzzy_space R)
          (phi0 : judgement sig (fcarrier X0)),
        E X0 phi0 -> F X0 phi0
H : E X phi
\end{gctx}
\tcblower
\begin{gconc}
|- F X phi
\end{gconc}
    \end{goalbox}
    \smallskip
    {\footnotesize The rule \rid{FD\_Init} turns ``derivable in
    \rid{F}'' back into ``an axiom of \rid{F}'', which is exactly what
    \rid{Hincl} supplies.}
  \end{column}
  \end{columns}
\end{frame}

\begin{frame}[fragile]{Lemma \rid{frel\_derives\_weakening}}
  \begin{columns}[T]
  \begin{column}{0.42\textwidth}
    \begin{scriptbox}\pstep{5}\weakeningscript\end{scriptbox}
  \end{column}
  \begin{column}{0.56\textwidth}
    \begin{goalbox}[title={Goal 1 / 12, after \texttt{apply: Hincl}}]
\begin{gctx}
R : realType
sig : signature
E, F : fuzzy_theory R sig
mode : frel_derivation_mode
X : fuzzy_space R
phi : judgement sig (fcarrier X)
H : E X phi
\end{gctx}
\tcblower
\begin{gconc}
|- E X phi
\end{gconc}
    \end{goalbox}
    \smallskip
    {\footnotesize \rid{Hincl} has disappeared from the context:
    ssreflect's \rid{apply:} consumes the hypothesis it applies. What
    is left is exactly the assumption \rid{H} handed to us by
    \rid{FD\_Init}.}
  \end{column}
  \end{columns}
\end{frame}

\begin{frame}[fragile]{Lemma \rid{frel\_derives\_weakening}}
  \begin{columns}[T]
  \begin{column}{0.42\textwidth}
    \begin{scriptbox}\pstep{6}\weakeningscript\end{scriptbox}
    \smallskip
    {\footnotesize The \rid{FD\_Init} case is closed. \textbf{11 goals}
    remain.}
  \end{column}
  \begin{column}{0.56\textwidth}
    \begin{goalbox}[title={New Goal 1 / 11 \quad (FD\_EqRefl)}]
\begin{gctx}
R : realType
sig : signature
E, F : fuzzy_theory R sig
Hincl : forall (X0 : fuzzy_space R)
          (phi : judgement sig (fcarrier X0)),
        E X0 phi -> F X0 phi
mode : frel_derivation_mode
X : fuzzy_space R
s : term sig (fcarrier X)
\end{gctx}
\tcblower
\begin{gconc}
|- frel_derives F mode (EqJ s s)
\end{gconc}
    \end{goalbox}
    \smallskip
    {\footnotesize From here on every goal has the same shape:
    re-apply the rule the case came from, to the induction
    hypotheses.}
  \end{column}
  \end{columns}
\end{frame}

\begin{frame}[fragile]{Lemma \rid{frel\_derives\_weakening}}
  \begin{columns}[T]
  \begin{column}{0.42\textwidth}
    \begin{scriptbox}\pstep{7}\weakeningscript\end{scriptbox}
  \end{column}
  \begin{column}{0.56\textwidth}
    \begin{goalbox}[title={Proof state}]
\begin{gconc}
No more goals.
\end{gconc}
    \end{goalbox}
    \smallskip
    \begin{alertblock}{Why the eleven remaining goals go by automation}
      \small
      Each constructor's premises mention \rid{frel\_derives E}, and the
      induction hypotheses turn every one of them into the corresponding
      \rid{frel\_derives F} fact. So \rid{eauto} only has to pick the
      matching constructor --- including \rid{FD\_OrderComplete}, whose
      premise is a \emph{function} over $\delta$.
    \end{alertblock}
  \end{column}
  \end{columns}
\end{frame}

% ============================================================
\section{Soundness}
% ============================================================

\begin{frame}{Soundness}
  Let $(A,d)$ be a fuzzy relation and $s,t$ two barycentric algebra
  terms with variables in $A$. Then
  \vfill
  \begin{block}{}
    If $\;\ICA \vdash \forall(A,d_A).\, s \qeq{\eps} t\;$ is derivable
    \vspace{0.5em}

    then $\;\Kan(d_A)(\den{s},\den{t}) \leq \eps$.
  \end{block}
  \vfill
  \pause
  \begin{alertblock}{$\vdash$ versus $\models$}
    \begin{description}\small
      \item[$E \vdash \phi$] \emph{syntactic}: there is a derivation of
        $\phi$ using the rules of the proof system and the axioms in
        $E$. In Rocq: \rid{frel\_derives E mode phi}.
      \item[$E \models \phi$] \emph{semantic}: every quantitative
        algebra $\mathbb{A}$ with $\mathbb{A} \models E$ also satisfies
        $\phi$. In Rocq: \rid{fuzzy\_models A E -> satisfies A X phi},
        universally quantified over \rid{A}.
    \end{description}
    \vspace{0.4em}
    Soundness is $\vdash\;\Rightarrow\;\models$; completeness is
    $\models\;\Rightarrow\;\vdash$.
  \end{alertblock}
\end{frame}

\newcommand{\sndpane}[1]{%
  \begin{scriptbox}[fontupper=\ttfamily\fontsize{5.8pt}{6.6pt}\selectfont\frenchspacing]%
  \pstep{#1}\soundnessscript\end{scriptbox}}

\begin{frame}[fragile]{Theorem \rid{derives\_fin\_sound}}
  \begin{columns}[T]
  \begin{column}{0.44\textwidth}\sndpane{1}\end{column}
  \begin{column}{0.54\textwidth}
    \begin{goalbox}[title={Proof state}]
\begin{gctx}
R : realType
sig : signature
A : FuzzyAlgebra R sig
E : fuzzy_theory R sig
X : fuzzy_space R
phi : judgement sig (fcarrier X)
\end{gctx}
\tcblower
\begin{gconc}
|- fuzzy_models A E ->
   derives_fin E phi ->
   satisfies A X phi
\end{gconc}
    \end{goalbox}
    \smallskip
    {\footnotesize This is $\vdash\;\Rightarrow\;\models$: the
    conclusion \rid{satisfies A X phi} unfolds to a claim about
    \emph{every} $1$-Lipschitz interpretation.}
  \end{column}
  \end{columns}
\end{frame}

\begin{frame}[fragile]{Theorem \rid{derives\_fin\_sound}}
  \begin{columns}[T]
  \begin{column}{0.44\textwidth}\sndpane{2}\end{column}
  \begin{column}{0.54\textwidth}
    \begin{goalbox}[title={Proof state}]
\begin{gctx}
Hmodel : forall (X0 : fuzzy_space R)
           (phi0 : judgement sig (fcarrier X0)),
         E X0 phi0 ->
         forall rho0 : interpretation X0 A,
         match phi0 with
         | EqJ s t => fuzzy_eval A rho0 s = fuzzy_eval A rho0 t
         | QEqJ eps s t =>
             frel (fa_space A) (fuzzy_eval A rho0 s)
                               (fuzzy_eval A rho0 t) <= eps
         end
Hderive : frel_derives E FRelFinite phi
rho : interpretation X A
\end{gctx}
\tcblower
\begin{gconc}
|- match phi with
   | EqJ s t => fuzzy_eval A rho s = fuzzy_eval A rho t
   | QEqJ eps s t =>
       frel (fa_space A) (fuzzy_eval A rho s)
                         (fuzzy_eval A rho t) <= eps
   end
\end{gconc}
    \end{goalbox}
  \end{column}
  \end{columns}
\end{frame}

\begin{frame}[fragile]{Theorem \rid{derives\_fin\_sound}}
  \begin{columns}[T]
  \begin{column}{0.44\textwidth}\sndpane{3}\end{column}
  \begin{column}{0.54\textwidth}
    \begin{goalbox}[title={Proof state}]
\begin{gctx}
Hmodel : ... (as before)
mode : frel_derivation_mode
Hderive : frel_derives E mode phi
rho : interpretation X A
\end{gctx}
\tcblower
\begin{gconc}
|- mode = FRelFinite ->
   match phi with ... end
\end{gconc}
    \end{goalbox}
    \smallskip
    \begin{alertblock}{The \rid{remember} trick}
      \small
      \rid{derives\_fin} pins \rid{mode := FRelFinite}, but
      \rid{induction} needs \rid{mode} \emph{general} --- otherwise
      Rocq refuses to generalize the index. So we abstract it and carry
      the equation \rid{Hmode} into the goal, to be used case by case.
    \end{alertblock}
  \end{column}
  \end{columns}
\end{frame}

\begin{frame}[fragile]{Theorem \rid{derives\_fin\_sound}}
  \begin{columns}[T]
  \begin{column}{0.44\textwidth}\sndpane{4}\end{column}
  \begin{column}{0.54\textwidth}
    \begin{goalbox}[title={Goal 1 / 12 \quad (FD\_Init)}]
\begin{gctx}
Hmodel : ... (as before)
mode : frel_derivation_mode
X : fuzzy_space R
phi : judgement sig (fcarrier X)
H : E X phi
rho : interpretation X A
Hmode : mode = FRelFinite
\end{gctx}
\tcblower
\begin{gconc}
|- match phi with ... end
\end{gconc}
    \end{goalbox}
    \smallskip
    {\footnotesize \textbf{12 goals}, one per rule. Here \rid{H} says
    \rid{phi} is an axiom of \rid{E}, and \rid{Hmodel} says \rid{A}
    satisfies every axiom of \rid{E}: one application closes it.}
  \end{column}
  \end{columns}
\end{frame}

\begin{frame}[fragile]{Theorem \rid{derives\_fin\_sound}}
  \begin{columns}[T]
  \begin{column}{0.44\textwidth}\sndpane{5}\end{column}
  \begin{column}{0.54\textwidth}
    \begin{goalbox}[title={Goal 5 / 12 \quad (FD\_EqCong)}]
\begin{gctx}
f : sym sig
s, t : 'I_(arity f) -> term sig (fcarrier X)
H : forall i, frel_derives E mode (EqJ (s i) (t i))
H0 : forall (i : 'I_(arity f)) (rho0 : interpretation X A),
     mode = FRelFinite ->
     fuzzy_eval A rho0 (s i) = fuzzy_eval A rho0 (t i)
rho : interpretation X A
Hmode : mode = FRelFinite
\end{gctx}
\tcblower
\begin{gconc}
|- fuzzy_eval A rho (App f s) = fuzzy_eval A rho (App f t)
\end{gconc}
    \end{goalbox}
    \smallskip
    {\footnotesize The equality cases are just \rid{symmetry} and
    \rid{rewrite} with the IHs. Here both sides reduce to
    \rid{fa\_ops A f} applied to the argument vectors, \rid{H0} equates
    those pointwise, and functional extensionality finishes.}
  \end{column}
  \end{columns}
\end{frame}

\begin{frame}[fragile]{Theorem \rid{derives\_fin\_sound}}
  \begin{columns}[T]
  \begin{column}{0.44\textwidth}\sndpane{6}\end{column}
  \begin{column}{0.54\textwidth}
    \begin{goalbox}[title={Goal 6 / 12 \quad (FD\_UseVariables)}]
\begin{gctx}
x, y : fcarrier X
rho : interpretation X A
Hmode : mode = FRelFinite
\end{gctx}
\tcblower
\begin{gconc}
|- frel (fa_space A) (fuzzy_eval A rho (Var x))
                     (fuzzy_eval A rho (Var y))
   <= frel X x y
\end{gconc}
    \end{goalbox}
    \smallskip
    {\footnotesize This is \emph{precisely} why an \rid{interpretation}
    carries its own $1$-Lipschitz proof. The axiom
    ``$x \qeq{d_B(x,y)} y$'' is sound exactly because $\den{\cdot}$ is
    non-expansive --- so the goal \emph{is} \rid{interpretation\_nexp}.}
  \end{column}
  \end{columns}
\end{frame}

\begin{frame}[fragile]{Theorem \rid{derives\_fin\_sound}}
  \begin{columns}[T]
  \begin{column}{0.44\textwidth}\sndpane{7}\end{column}
  \begin{column}{0.54\textwidth}
    \begin{goalbox}[title={Goal 7 / 12 \quad (FD\_Subst)}]
\begin{gctx}
sigma : fcarrier X -> term sig (fcarrier Y)
H0 : forall (x y : fcarrier X) (rho0 : interpretation Y A),
     mode = FRelFinite ->
     frel (fa_space A) (fuzzy_eval A rho0 (sigma x))
                       (fuzzy_eval A rho0 (sigma y))
     <= frel X x y
IHfrel_derives : forall rho0 : interpretation X A,
     mode = FRelFinite -> match phi with ... end
rho : interpretation Y A
\end{gctx}
\tcblower
\begin{gconc}
|- match subst_judgement sigma phi with ... end
\end{gconc}
    \end{goalbox}
    {\footnotesize \rid{H0} says exactly that
    $z \mapsto \den{\sigma(z)}$ is $1$-Lipschitz --- enough to
    \emph{build} an interpretation and feed it to the IH:}
    \begin{srcbox}
  pose rho_sigma := @Build_interpretation R sig X A
    (fun z => fuzzy_eval A rho (sigma z)) Hnexp.
    \end{srcbox}
  \end{column}
  \end{columns}
\end{frame}

\begin{frame}[fragile]{Theorem \rid{derives\_fin\_sound}}
  \begin{columns}[T]
  \begin{column}{0.44\textwidth}\sndpane{8}\end{column}
  \begin{column}{0.54\textwidth}
    \begin{goalbox}[title={Goal 12 / 12 \quad (FD\_OrderComplete)}]
\begin{gctx}
eps : R
s, t : term sig (fcarrier X)
H  : forall delta : R, eps < delta ->
     frel_derives E FRelFull (QEqJ delta s t)
H0 : forall delta : R, eps < delta ->
     forall rho0 : interpretation X A,
     FRelFull = FRelFinite -> ...
rho : interpretation X A
Hmode : FRelFull = FRelFinite
\end{gctx}
\tcblower
\begin{gconc}
|- frel (fa_space A) (fuzzy_eval A rho s)
                     (fuzzy_eval A rho t) <= eps
\end{gconc}
    \end{goalbox}
    \smallskip
    {\footnotesize The infinitary rule \alert{cannot have been used}: it
    only produces derivations at \rid{FRelFull}. The \rid{remember}
    trick turned that observation into \rid{discriminate Hmode}.}
  \end{column}
  \end{columns}
\end{frame}

\begin{frame}[fragile]{The same theorem at \rid{FRelFull}}
  \rid{derives\_full\_sound} has an identical script except for the
  last case, where \rid{Hmode : mode = FRelFull} is now
  \emph{consistent}:
  \begin{goalbox}[title={\texttt{derives\_full\_sound}, Goal 12/12}]
\begin{gctx}
H0 : forall delta : R, eps < delta -> forall rho0 : interpretation X A,
     mode = FRelFull ->
     frel (fa_space A) (fuzzy_eval A rho0 s) (fuzzy_eval A rho0 t) <= delta
Hmode : mode = FRelFull
\end{gctx}
\tcblower
\begin{gconc}
|- frel (fa_space A) (fuzzy_eval A rho s) (fuzzy_eval A rho t) <= eps
\end{gconc}
  \end{goalbox}
  \begin{srcbox}
Lemma le_of_all_strict_upper_bounds {R : realType} (a eps : R) :
  (forall delta, (eps < delta)%R -> (a <= delta)%R) -> (a <= eps)%R.
Proof.
  move => H. apply /Num.Theory.ler_addgt0Pr => e He.
  apply: H. by rewrite Num.Theory.ltrDl.
Qed.
  \end{srcbox}
  \vfill
  So the infinitary rule is sound too --- it is \emph{eliminability},
  not soundness, that is the hard question.
\end{frame}

\begin{frame}{Completeness}
  Let $(A,d)$ be a fuzzy relation and $s,t$ two barycentric algebra
  terms with variables in $A$. Then
  \vfill
  \begin{block}{}
    If $\;\Kan(d_A)(\den{s},\den{t}) \leq \eps$
    \vspace{0.5em}

    then $\;\ICA \vdash \forall(A,d_A).\, s \qeq{\eps} t\;$ is
    derivable.
  \end{block}
  \vfill
  \pause
  {\small And --- this is the compactness result --- it is derivable
  by a finite proof, i.e.\ already in
  \rid{derives\_fin}.}
\end{frame}

% ============================================================
\section{Instantiating: $\ICA$ and the Kantorovich model}
% ============================================================

\begin{frame}[fragile]{Formalization: the $\ICA$ signature}
  \file{src/ICA.v}
  \begin{srcbox}
Record ica_weight (R : realType) := {
  ica_weight_val : R;
  ica_weight_open : (0 < ica_weight_val < 1)%R;
}.

Inductive ica_sym (R : realType) : Type := | ica_plus : ica_weight R -> ica_sym R.
Definition ica_arity {R : realType} (_ : ica_sym R) : nat := 2.
Definition ica_signature {R : realType} : signature :=
  {| sym := ica_sym R; arity := ica_arity |}.

Definition ica_op {R : realType} {X : Type} (p : ica_weight R)
    (x y : term ica_signature X) : term ica_signature X :=
  @App ica_signature X (ica_plus p)
    (fun i => if (Nat.eqb (i : nat) 0) then x else y).

Notation "x <+ p +> y" := (ica_op p x y).
  \end{srcbox}
  \vfill
  The signature is indexed by the open unit interval of an
  abstract \rid{realType} --- genuinely uncountable, which is why
  \rid{sym} had to be a \rid{Type} rather than a \rid{finType}.
\end{frame}

\begin{frame}[fragile]{Formalization: weights and contexts}
  The weights appearing in the axioms are themselves constructions:
  \begin{srcbox}
Definition ica_weight_complement (p : ica_weight R) : ica_weight R.
  (* 1 - p, openness via [subr_gt0] / [gtrBl] *)
Definition ica_weight_product (p q : ica_weight R) : ica_weight R.
  (* p * q, via [mulr_gt0] and [mulr_ilt1] *)
Definition ica_weight_assoc_inner (p q : ica_weight R) : ica_weight R.
  (* (1-p)q / (1-pq) -- the skew-associativity weight *)
  \end{srcbox}
  and so are the fuzzy relations $d_B$ used as contexts:
  \begin{srcbox}
Definition ica_full_space (R : realType) (n : nat) : fuzzy_space R.
  (* 'I_n with frel = fun _ _ => 1 : the discrete fuzzy relation *)

(*  [ 1, 1, eps, 1     ]     [ 1, 1, 1,   1 ]                  *)
(*  [ 1, 1, 1,   delta ]     [ 1, 1, 1,   1 ]                  *)
Definition ica_interp_rel {R : realType} (eps delta : R) (a b : 'I_4) : R :=
  match nat_of_ord a, nat_of_ord b with
  | 0, 2 => eps  | 1, 3 => delta  | _, _ => 1
  end.
  \end{srcbox}
  The matrix drawn on the ``interpolative'' slide, verbatim.
\end{frame}

\begin{frame}[fragile]{Formalization: \rid{Inductive ica\_theory}}
  \begin{srcbox}
Inductive ica_theory (R : realType) : fuzzy_theory R (@ica_signature R) :=
  | ICA_Idem (p : ica_weight R) :
      @ica_theory R (ica_full_space R 1)
        (EqJ ((Var (inord 0)) <+ p +> (Var (inord 0))) (Var (inord 0)))
  | ICA_Skew_Comm (p : ica_weight R) :
      @ica_theory R (ica_full_space R 2)
        (EqJ ((Var (inord 0)) <+ p +> (Var (inord 1)))
             ((Var (inord 1)) <+ ica_weight_complement p +> (Var (inord 0))))
  | ICA_Skew_Assoc (p q : ica_weight R) :
      @ica_theory R (ica_full_space R 3)
        (EqJ (((Var (inord 0)) <+ p +> (Var (inord 1))) <+ q +> (Var (inord 2)))
             ((Var (inord 0)) <+ ica_weight_product p q +>
               ((Var (inord 1)) <+ ica_weight_assoc_inner p q +> (Var (inord 2)))))
  | ICA_Interp (p : ica_weight R) (eps delta : R)
      (Heps : (0 <= eps <= 1)%R) (Hdelta : (0 <= delta <= 1)%R) :
      @ica_theory R (@ica_interp_space R eps delta Heps Hdelta)
        (QEqJ (weighted_sum (ica_probability_weight p) eps delta)
          ((Var (inord 0)) <+ p +> (Var (inord 1)))
          ((Var (inord 2)) <+ p +> (Var (inord 3)))).
  \end{srcbox}
  \vfill
  An \rid{Inductive} \emph{in} \rid{fuzzy\_theory R sig}: each
  constructor picks its own context \rid{X}, so the four axiom
  \emph{schemes} become four constructors, each quantifying over its
  parameters $p,q,\eps,\delta$.
\end{frame}

\begin{frame}[fragile]{Formalization: instantiating the axioms}
  Using an axiom requires the substitution rule, and its
  ``provably $1$-Lipschitz'' premise --- discharged by \rid{FD\_Max}
  because the context is discrete:
  \begin{srcbox}
Lemma ica_idem_instance {R : realType} (mode : frel_derivation_mode)
    (X : fuzzy_space R) (p : ica_weight R)
    (t : term (@ica_signature R) (fcarrier X)) :
  @frel_derives R ica_signature (@ica_theory R) mode X (EqJ (t <+ p +> t) t).
Proof.
  (* ... *)
  apply: (FD_Subst (X := ica_full_space R 1) (sigma := fun _ => t)).
  - apply: FD_Init. exact: ICA_Idem p.
  - move => x y. apply: FD_Max.       (* frel (ica_full_space R 1) x y = 1 *)
Qed.
  \end{srcbox}
  \vfill
  Similarly \rid{ica\_skew\_comm\_instance},
  \rid{ica\_skew\_assoc\_instance} and \rid{ica\_plus\_congr} in
  \file{src/BarycentricNormalization.v}: they turn the axioms into
  ordinary rewriting steps on arbitrary terms.
\end{frame}

\begin{frame}[fragile]{Formalization: distributions and couplings}
  \file{src/ProbabilityDistribution.v}
  \begin{srcbox}
Record probability_distribution (R : realType) (X : finType) := {
  probability_mass :> X -> R;
  probability_mass_ge0 : forall x, (0 <= probability_mass x)%R;
  probability_mass_total : \sum_(x : X) probability_mass x = 1;
}.
Record coupling {R : realType} {X Y : finType}
    (mu : probability_distribution R X) (nu : probability_distribution R Y) := {
  coupling_distribution : joint_probability_distribution R X Y;
  coupling_fst : forall x, \sum_(y : Y) coupling_distribution x y = mu x;
  coupling_snd : forall y, \sum_(x : X) coupling_distribution x y = nu y;
}.
Definition coupling_cost {R : realType} {X Y : finType} (d : X -> Y -> R)
    {mu} {nu} (gamma : coupling mu nu) : R :=
  \sum_(x : X) \sum_(y : Y) coupling_distribution gamma x y * d x y.

Definition kantorovich_lifting {R : realType} {X : finType} (d : X -> X -> R)
    (mu nu : probability_distribution R X) : R :=
  inf [set r : R | exists gamma : coupling mu nu, r = coupling_cost d gamma].
  \end{srcbox}
  \rid{X : finType} makes ``finitely supported'' automatic; \rid{inf}
  is MathComp-Analysis' infimum on \rid{set R}.
\end{frame}

\begin{frame}[fragile]{Formalization: the facts the model needs}
  \begin{srcbox}
Definition convex_mixture (p : probability_weight R)
    (mu nu : probability_distribution R X) : probability_distribution R X.
  (* pointwise [weighted_sum p (mu x) (nu x)]; totality from
     [sum_weighted_sum] and [weighted_sum_one] *)

Lemma coupling_convex_mixture (p : probability_weight R)
    (gamma1 : coupling mu1 nu1) (gamma2 : coupling mu2 nu2) :
  coupling (convex_mixture p mu1 mu2) (convex_mixture p nu1 nu2).

Lemma fuzzy_kantorovich_lifting (d : X -> X -> R) (mu nu : ...) :
  (forall x y, 0 <= d x y <= 1) -> 0 <= kantorovich_lifting d mu nu <= 1.
  \end{srcbox}
  Lemma 2.5 of the paper: a convex combination of couplings is a
  coupling of the convex combinations --- the witness the interpolative
  axiom needs.
\end{frame}

\begin{frame}[fragile]{Formalization: convexity of the lifting}
  \file{src/KantorovichProperties.v}
  \begin{srcbox}
Lemma kantorovich_almost_optimal (d : X -> X -> R) (mu nu : ...) (eta : R) :
  (forall x y, (0 <= d x y)%R) -> (0 < eta)%R ->
  exists gamma : coupling mu nu,
    (coupling_cost d gamma < kantorovich_lifting d mu nu + eta)%R.

Lemma kantorovich_convex_mixture (d : X -> X -> R) (p : probability_weight R) :
  (forall x y, (0 <= d x y <= 1)%R) ->
  (kantorovich_lifting d (convex_mixture p mu1 mu2) (convex_mixture p nu1 nu2)
     <= weighted_sum p (kantorovich_lifting d mu1 nu1)
                       (kantorovich_lifting d mu2 nu2))%R.
  \end{srcbox}
  \vfill
  \begin{alertblock}{}
    No compactness of $\Cpl(\mu,\nu)$ is needed here:
    \emph{almost}-optimal couplings suffice, obtained from
    \rid{inf\_adherent} and glued with \rid{ler\_addgt0Pr}. Compactness
    is only needed for \emph{finite} completeness proofs.
  \end{alertblock}
\end{frame}

\begin{frame}[fragile]{Formalization: the Kantorovich model}
  \file{src/ICAKantorovichSoundness.v}
  \begin{srcbox}
Section KantorovichModel.
Context {R : realType} {X : finType}.
Variable d : X -> X -> R.
Hypothesis Hd : forall x y, (0 <= d x y <= 1)%R.

Definition finite_fuzzy_space : fuzzy_space R :=
  {| fcarrier := X; frel := d; frel_range := Hd |}.

Definition kantorovich_fuzzy_space : fuzzy_space R :=
  {| fcarrier := probability_distribution R X;
     frel := kantorovich_lifting d;
     frel_range := fun mu nu => fuzzy_kantorovich_lifting mu nu Hd |}.

Definition kantorovich_algebra : FuzzyAlgebra R (@ica_signature R) :=
  {| fa_space := kantorovich_fuzzy_space; fa_ops := kantorovich_ops |}.
  \end{srcbox}
  \vfill
  \rid{kantorovich\_ops} interprets \rid{ica\_plus p} as
  \rid{convex\_mixture}. So this record \emph{is}
  $(\Dist(A),\{\bpl{p}\},\Kan(d))$.
\end{frame}

\newcommand{\modpane}[1]{%
  \begin{scriptbox}[fontupper=\ttfamily\fontsize{5.8pt}{6.6pt}\selectfont\frenchspacing]%
  \pstep{#1}\modelscript\end{scriptbox}}

\begin{frame}[fragile]{Modelhood: $(\Dist(A),\{\bpl{p}\},\Kan(d)) \models \ICA$}
  \begin{columns}[T]
  \begin{column}{0.44\textwidth}\modpane{1}\end{column}
  \begin{column}{0.54\textwidth}
    \begin{goalbox}[title={Goal 1 / 4 \quad (ICA\_Idem)}]
\begin{gctx}
d : X -> X -> R
Hd : forall x y : X, 0 <= d x y <= 1
p : ica_weight R
rho : interpretation (ica_full_space R 1)
        kantorovich_algebra
\end{gctx}
\tcblower
\begin{gconc}
|- fuzzy_eval kantorovich_algebra rho
     (Var (inord 0) <+ p +> Var (inord 0)) =
   fuzzy_eval kantorovich_algebra rho (Var (inord 0))
\end{gconc}
    \end{goalbox}
    \smallskip
    {\footnotesize \rid{case: Hphi} on the \emph{inductive} theory gives
    one goal per axiom --- four in total. The first three are the
    convex-algebra equations, the fourth is interpolation.}
  \end{column}
  \end{columns}
\end{frame}

\begin{frame}[fragile]{Modelhood: $(\Dist(A),\{\bpl{p}\},\Kan(d)) \models \ICA$}
  \begin{columns}[T]
  \begin{column}{0.44\textwidth}\modpane{2}\end{column}
  \begin{column}{0.54\textwidth}
    \begin{goalbox}[title={Goal 4 / 4 \quad (ICA\_Interp)}]
\begin{gctx}
p : ica_weight R
eps, delta : R
Heps : 0 <= eps <= 1
Hdelta : 0 <= delta <= 1
rho : interpretation (ica_interp_space Heps Hdelta)
        kantorovich_algebra
\end{gctx}
\tcblower
\begin{gconc}
|- frel (fa_space kantorovich_algebra)
     (fuzzy_eval kantorovich_algebra rho
        (Var (inord 0) <+ p +> Var (inord 1)))
     (fuzzy_eval kantorovich_algebra rho
        (Var (inord 2) <+ p +> Var (inord 3)))
   <= weighted_sum (ica_probability_weight p) eps delta
\end{gconc}
    \end{goalbox}
    \smallskip
    {\footnotesize The three equational axioms fall to the
    \rid{convex\_mixture\_*} lemmas: $\Dist(X)$ \emph{is} a convex
    algebra, so they are equalities of distributions, checked
    pointwise. Only interpolation says anything quantitative.}
  \end{column}
  \end{columns}
\end{frame}

\begin{frame}[fragile]{Modelhood: $(\Dist(A),\{\bpl{p}\},\Kan(d)) \models \ICA$}
  \begin{columns}[T]
  \begin{column}{0.44\textwidth}\modpane{3}\end{column}
  \begin{column}{0.54\textwidth}
    \begin{goalbox}[title={Goal 4 / 4, after \texttt{rewrite !kantorovich\_eval\_op}}]
\begin{gconc}
|- frel (fa_space kantorovich_algebra)
     (convex_mixture (ica_probability_weight p)
        (fuzzy_eval kantorovich_algebra rho (Var (inord 0)))
        (fuzzy_eval kantorovich_algebra rho (Var (inord 1))))
     (convex_mixture (ica_probability_weight p)
        (fuzzy_eval kantorovich_algebra rho (Var (inord 2)))
        (fuzzy_eval kantorovich_algebra rho (Var (inord 3))))
   <= weighted_sum (ica_probability_weight p) eps delta
\end{gconc}
    \end{goalbox}
    \smallskip
    {\footnotesize Now the goal has exactly the shape of
    \rid{kantorovich\_convex\_mixture}:
    \[
      \Kan(d)(\mu_1 \bpl{p} \mu_2,\; \nu_1 \bpl{p} \nu_2)
      \;\leq\; p\,\Kan(d)(\mu_1,\nu_1) + (1-p)\,\Kan(d)(\mu_2,\nu_2).
    \]}
  \end{column}
  \end{columns}
\end{frame}

\begin{frame}[fragile]{Modelhood: $(\Dist(A),\{\bpl{p}\},\Kan(d)) \models \ICA$}
  \begin{columns}[T]
  \begin{column}{0.44\textwidth}\modpane{4}\end{column}
  \begin{column}{0.54\textwidth}
    \begin{goalbox}[title={Goal 1 / 2, after \texttt{apply: weighted\_sum\_le}}]
\begin{gconc}
|- kantorovich_lifting d
     (fuzzy_eval kantorovich_algebra rho (Var (inord 0)))
     (fuzzy_eval kantorovich_algebra rho (Var (inord 2)))
   <= eps
\end{gconc}
    \end{goalbox}
    \begin{goalbox}[title={Goal 2 / 2}]
\begin{gconc}
|- kantorovich_lifting d
     (fuzzy_eval kantorovich_algebra rho (Var (inord 1)))
     (fuzzy_eval kantorovich_algebra rho (Var (inord 3)))
   <= delta
\end{gconc}
    \end{goalbox}
    {\footnotesize These two are \emph{exactly}
    $d(\den{x},\den{w}) \leq \eps$ and
    $d(\den{y},\den{z}) \leq \delta$ --- the $1$-Lipschitz property of
    \rid{rho}, read off the context matrix at
    $(0,2) \mapsto \eps$ and $(1,3) \mapsto \delta$. The paper's proof,
    in three tactics.}
  \end{column}
  \end{columns}
\end{frame}

\begin{frame}[fragile]{Formalization: soundness, instantiated (Corollary 4.7)}
  \begin{srcbox}
Definition dirac_interpretation :
    interpretation finite_fuzzy_space kantorovich_algebra :=
  @Build_interpretation R (@ica_signature R)
    finite_fuzzy_space kantorovich_algebra
    (fun x : X => dirac x)
    (fun x y => kantorovich_dirac_le x y Hd).   (* Proposition 4.5 *)

Definition ica_term_distribution (s : term (@ica_signature R) X) :=
  fuzzy_eval kantorovich_algebra (fun x : X => dirac x) s.

Theorem finite_kantorovich_soundness (s t : term (@ica_signature R) X) (eps : R) :
  @derives_full R (@ica_signature R) (@ica_theory R)
    finite_fuzzy_space (QEqJ eps s t) ->
  (kantorovich_lifting d (ica_term_distribution s)
                         (ica_term_distribution t) <= eps)%R.
Proof.
  move => H.
  apply: (derives_full_sound kantorovich_models_ica H dirac_interpretation).
Qed.
  \end{srcbox}
  Soundness proper is \rid{derives\_full\_sound}; Dirac is not a
  restriction, it is the one instance whose values are the
  distributions $\den{s},\den{t}$ denoted by the terms.
\end{frame}

% ============================================================
\section{Where the development stands}
% ============================================================

\begin{frame}[fragile]{Current state of completeness}
  Done: \file{Signature}, \file{FuzzyRelation}, \file{FRelDeduction}
  (both modes sound), \file{ProbabilityDistribution},
  \file{KantorovichProperties}, \file{ICA},
  \file{ICAKantorovichSoundness}. Plus \file{QET}/\file{Metric}/%
  \file{Category} for the metric view.
  \medskip

  In progress --- \file{src/BarycentricNormalization.v}:
  \begin{srcbox}
Fixpoint term_distribution (s : term (@ica_signature R) X)
  : probability_distribution R X.       (* the denotation [s] *)
Fixpoint term_vars (s : term (@ica_signature R) X) : {set X}.

Lemma term_distribution_support (s : term (@ica_signature R) X) :
  distribution_support (term_distribution s) = term_vars s.

Lemma condition_recover (mu : probability_distribution R X) (x : X) :
  (mu x < 1) -> x \in distribution_support mu ->
  convex_mixture (ica_probability_weight (ica_weight_clamp (mu x)))
    (dirac x) (condition mu x) = mu.
  \end{srcbox}
  \rid{condition\_recover} is the peeling step: any $\mu$ is
  $\delta_x \bpl{\mu(x)} (\mu \mid {\neq} x)$, and the remainder has
  strictly smaller support.
\end{frame}

\begin{frame}{How to proceed}
  \begin{enumerate}
    \item \textbf{Normal forms.} Iterate \rid{condition\_recover} to
      write every term as an $n$-ary convex combination
      $\sum_i p_i x_i$ over its support --- the induction that
      \rid{term\_distribution\_support} makes well-founded.
      \medskip
    \item \textbf{Lemma 4.9.} For $e(\bm x) = \sum_i p_i x_i$ and
      tuples $\bm a, \bm b$:
      \[
        \ICA \vdash \forall(A,d).\, e(\bm a) \qeq{\eps} e(\bm b),
        \qquad \eps = \textstyle\sum_i p_i\, d(a_i,b_i),
      \]
      by induction on $n$, using one instance of the interpolative
      axiom per step. Purely syntactic, and finite.
      \medskip
    \item \textbf{Proposition 4.10.} A coupling $\gamma \in
      \Cpl(\den{s},\den{t})$ is a recipe for rewriting $s$ and $t$ into
      the \emph{same shape} $e(\bm x)$; then Lemma 4.9 gives
      $s \qeq{V_d(\gamma)} t$, and \rid{FD\_Up} weakens to $\eps$.
      \medskip
    \item \textbf{Finiteness.} Step 3 needs \emph{one} coupling, not the
      infimum: with an optimal coupling the proof is finite, giving
      compactness of $\ICA$. \\
      {\small $\hookrightarrow$ requires compactness of
      $\Cpl(\mu,\nu)$ --- the one genuinely topological ingredient
      still missing.}
  \end{enumerate}
\end{frame}

% ============================================================
\section{Outlook: the quantum side}
% ============================================================

\begin{frame}{How does this relate to quantum?}
  \begin{itemize}
    \item[--] In my PhD I am trying to apply this line of thinking to
      quantum computing.
      \medskip
    \item[--] We want to start small, by finding a theory for
      mixed quantum states.
      \medskip
    \item[--] Later we would like a theory for programs.
  \end{itemize}
\end{frame}

\begin{frame}{Why mixed states?}
  \begin{itemize}
    \item[--] They can be viewed as probability distributions over pure
      states.
      \medskip
    \item[--] They are a model of the non-quantitative barycentric
      axioms: $\rho \bpl{p} \sigma = p\rho + (1-p)\sigma$ is again a
      density matrix.
      \medskip
    \item[--] There exists a notion of quantum optimal
      transport:
      \[
        T_C(\rho,\sigma) \;=\;
        \inf_{\Gamma \in \mathcal{C}(\rho,\sigma)} \operatorname{tr}(C\Gamma)
      \]
      where $\mathcal{C}(\rho,\sigma)$ are the quantum couplings ---
      bipartite states with marginals $\rho$ and $\sigma$.
  \end{itemize}
\end{frame}

\begin{frame}{Should be easy, right?}
  \begin{itemize}
    \item[--] Turns out that the barycentric structure can only reveal
      the transportation of mass.
      \medskip
    \item[--] Mixed states can be coupled in many different ways that
      transport the mass the same, but $\operatorname{tr}(C\Gamma)$
      gives different values.
      \medskip
      {\small $\hookrightarrow$ so the interpolative axiom, which is
      stated purely in terms of $\bpl{p}$, cannot separate them: the
      completeness argument of Proposition 4.10 breaks exactly at
      ``a coupling is a recipe for rewriting''.}
      \medskip
    \item[--] Currently using a mix of graphs and linear
      algebra, which seems promising.
  \end{itemize}
\end{frame}

% ============================================================
% ============================================================

\begin{frame}{Conclusion}
  \begin{itemize}
    \item[--] Quantitative equations $s \qeq{\eps} t$ replace exact
      equality by bounded distance; over \emph{fuzzy relations}
      the two judgement forms $=$ and $\qeq{\eps}$ must be kept apart.
      \medskip
    \item[--] Formalized in Rocq/MathComp: signatures and terms,
      fuzzy algebras, interpretations, the deduction system in
      two modes (finitary and order-complete) inside a single
      inductive.
      \medskip
    \item[--] Soundness is proved for both modes
      (\rid{derives\_fin\_sound}, \rid{derives\_full\_sound}), and the
      Kantorovich lifting on $\Dist(X)$ is proved to be a
      model of $\ICA$ (\rid{kantorovich\_models\_ica}), giving
      soundness of $\ICA$ with respect to the Kantorovich distance.
      \medskip
    \item[--] Completeness is under way: the normalization machinery
      (\rid{condition\_recover}, \rid{term\_distribution\_support}) is
      in place; Lemma 4.9 and Proposition 4.10 are next.
      \medskip
    \item[--] Goal: the same story for mixed quantum states,
      where the barycentric structure alone is not enough.
  \end{itemize}
  \vfill
  \begin{center}\large Thank you --- questions?\end{center}
\end{frame}

\appendix

\begin{frame}{References}
  \small
  \begin{itemize}
    \item R.~Mardare, P.~Panangaden, G.~D.~Plotkin.
      \emph{Quantitative algebraic reasoning}. LICS'16.
    \item R.~Mardare, P.~Panangaden, G.~D.~Plotkin.
      \emph{On the axiomatizability of quantitative algebras}. LICS'17.
    \item M.~Mio, R.~Sarkis, V.~Vignudelli.
      \emph{Universal quantitative algebra for fuzzy relations and
      generalised metric spaces}. LMCS 20(4:19), 2024.
    \item M.~Mio. \emph{Compact quantitative theories of convex
      algebras}. ENTICS, MFPS 2025.
    \item C.~Villani. \emph{Optimal Transport: Old and New}. Springer,
      2009.
    \item G.~Bacci, R.~Mardare, P.~Panangaden, G.~D.~Plotkin.
      \emph{Sum and tensor of quantitative effects}. LMCS 20(4:9), 2024.
  \end{itemize}
\end{frame}

\end{document}
